El backend expone su funcionalidad mediante una \acs{api} \acs{rest} que
sigue convenciones uniformes en todos sus recursos.

\begin{itemize}

  \item \textbf{Versionado.} Todos los \textit{endpoints} se exponen bajo
  el prefijo \texttt{/api/v1/}. Esta versión constituye un contrato
  estable para el \dfn{frontend}; cambios incompatibles requerirán un
  nuevo prefijo \texttt{/api/v2/}.

  \item \textbf{Convención de errores.} Todos los errores se devuelven
  con un formato uniforme que incluye un \texttt{error\_code} simbólico
  (no traducible en el \dfn{backend}) y, en el caso de errores de
  validación, un mapa por campo con los códigos correspondientes. El
  \dfn{frontend} traduce estos códigos a mensajes localizables.

  \item \textbf{Paginación, filtrado y ordenación.} Todas las listas se
  paginan mediante una paginación \textit{offset} (\texttt{page},
  \texttt{page\_size}) que devuelve el total de elementos y de páginas,
  lo que permite al \dfn{frontend} mostrar paginación clásica con número
  de página. El filtrado y la ordenación se aplican mediante parámetros
  de \textit{query} coherentes en todos los \textit{endpoints} listables.

  \item \textbf{Limitación de tasa.} El sistema aplica
  \textit{throttling} configurable por categoría de usuario (anónimo,
  autenticado estándar, gestor), con límites estrictos en
  \textit{endpoints} sensibles como el inicio de sesión.

  \item \textbf{Generación automática del contrato OpenAPI.} El esquema
  OpenAPI del SCDEE se deriva automáticamente del código mediante
  drf-spectacular, lo que garantiza que la documentación de la \acs{api}
  permanezca sincronizada con la implementación. El esquema completo,
  con la lista de \textit{endpoints} y sus contratos de entrada y salida,
  se incluye en \cref{AP:OPENAPI}.

\end{itemize}
